  \subsubsection{丸め誤差と係数の関係からガウス記号を外す}\label{節：ガウス型の丸め誤差}
    \,$a_{n+1}=\lfloor\lambda a_n\rfloor$では,整数との和の場合と違い,\,$\lambda$の整数部分を取り出すだけでは解けない.
    代わりに,
    \[
      \varepsilon_n=\lambda a_n-a_{n+1},
      \qquad 0\le\varepsilon_n<1
    \]
    とおき,丸めで捨てた部分を残して計算しよう.
    \textbf{誤差そのものではなく,誤差の範囲から次の整数部分を確定する}のが狙いである.

    係数$\lambda$が2次方程式を満たす場合は,\,$\lambda^2$を$\lambda$と整数で表せることが手掛かりになる.
    ただし,最後に現れる誤差がどの整数の間にあるかを示す必要がある.
    どの無理数の係数でも同じ隣接3項間漸化式へ帰着できるわけではない.

    \noindent \bookblocklink{example-gauss-role-5}{problem}{answer}{【例題】}
    \begin{enumerate}[label=(\arabic*),parsep=3pt,beginpenalty=10000,format=\bookitemlink{example-gauss-role-5}{problem}{answer}]
      \item 数列を次のように定める.ガウス記号を含まない漸化式に帰着させ,一般項を求めよ.
            \[
              a_1=1,\qquad
              a_{n+1}=\lfloor(1+\sqrt2)a_n\rfloor
              \quad(n\ge1).
            \]
    \end{enumerate}

    \noindent \bookblocklink{example-gauss-role-5}{hint}{problem}{【方針】}
    \begin{enumerate}[label=(\arabic*),parsep=3pt,beginpenalty=10000,format=\bookitemlink{example-gauss-role-5}{hint}{problem}]
      \item \,$\lambda=1+\sqrt2$とおき,\,$\lambda^2=2\lambda+1$を使います.
            整数から正の1未満の誤差を引く形にして,床関数の値を確定しましょう.
    \end{enumerate}

    \noindent \bookblocklink{example-gauss-role-5}{answer}{problem}{【解答】}
    \begin{enumerate}[label=(\arabic*),parsep=3pt,beginpenalty=10000,format=\bookitemlink{example-gauss-role-5}{answer}{problem}]
      \item \[
              a_1=1,\quad a_2=2,\qquad
              a_{n+2}=2a_{n+1}+a_n-1.
            \]
            \,$\alpha=1+\sqrt2,\ \beta=1-\sqrt2$とおくと,
            \[
              a_n=\frac12+\frac{\alpha^n+\beta^n}{4}.
            \]
    \end{enumerate}

    \noindent \bookblocklink{example-gauss-role-5}{solution}{problem}{【解法】}
    \begin{enumerate}[label=(\arabic*),parsep=3pt,beginpenalty=10000,format=\bookitemlink{example-gauss-role-5}{solution}{problem}]
      \item \,$\lambda=1+\sqrt2>2$とおく.正整数$a_n$から次の項も正整数となる.
            また,\,$\lambda a_n$は無理数なので,
            \[
              \varepsilon_n=\lambda a_n-a_{n+1}
            \]
            は$0<\varepsilon_n<1$を満たす.
            \,$\lambda^2=2\lambda+1$を使うと,
            \begin{align*}
              \lambda a_{n+1}
                &=\lambda(\lambda a_n-\varepsilon_n)\\
                &=2a_{n+1}+a_n-(\lambda-2)\varepsilon_n.
            \end{align*}
            \,$0<\lambda-2=\sqrt2-1<1$なので,最後に引く数も0と1の間にある.
            \,$M_n=2a_{n+1}+a_n$は整数だから,
            \[
              M_n-1<\lambda a_{n+1}<M_n.
            \]
            従って$a_{n+2}=M_n-1$となり,【解答】の隣接3項間漸化式を得る.
            初期条件は$a_2=\lfloor1+\sqrt2\rfloor=2$である.

            定数項を消すため$b_n=a_n-\frac12$とおくと,
            \[
              b_{n+2}=2b_{n+1}+b_n,
              \qquad b_1=\frac12,\quad b_2=\frac32.
            \]
            \sectionref{節：隣接3項間漸化式型}で学んだ特性方程式は$x^2-2x-1=0$で,
            その解が$\alpha,\beta$である.
            \,$\alpha+\beta=2,\ \alpha^2+\beta^2=6$より,
            \[
              b_n=\frac{\alpha^n+\beta^n}{4}
            \]
            は初期条件と漸化式を満たす.
            二つの初期値から各項は一意に決まるので,この式が一般項である.
            丸め誤差を評価した後は,既習の線形の漸化式に戻ることができた.
    \end{enumerate}

